\chapter*{Заключение}
\addcontentsline{toc}{chapter}{Заключение}

Несоответствие заключается в том, что оценки доступности и обеспеченности по нормативам отвечают на вопрос, где и что должно быть размещено, а ПКРТИ отвечает на вопрос, как развивать дорожно-транспортную инфраструктуру. В случае, если эти решения готовятся последовательно, возникает риск создания неэффективного плана застройки территории объектами обслуживания населения, поскольку согласование транспортного плана будет происходить на последующих этапах. Доступность определяется как совместная характеристика застройки и транспорта.

Таким образом, в существующих подходах транспортная доступность связывается с тремя изменяемыми компонентами -- объектами обслуживания, транспортной сетью и пространственным распределением спроса. Если маршрутная сеть проектируется с учетом пространственной доступности и географического равенства, то размещение объектов и проектирование маршрутов становятся частями одной задачи обеспечения доступности.

Общий результат исследований: доступность должна оцениваться по фактически допустимому состоянию сети; внешняя среда должна воздействовать на ребра через их класс, режим и качество; а итоговый КПЭ должен проверять достижимость допустимого пункта назначения, а не только общую топологическую связность графа. Таким образом, метод оценки сезонной устойчивости не сводится к одной интегральной метрике.

Объединяются изменения транспортной связности и размещение дополнительных сервисов, которые ранее рассматривались отдельно. Проверяемая гипотеза: улучшение связности между районами не более чем на 40\% может уменьшить минимальное число новых поликлиник, которые необходимо оптимально разместить \cite{kontsevik2024connectivityPlacement}. Маршрутное проектирование рассматривается как Transit Network Design Problem (TNDP): требуется определить допустимый набор линий общественного транспорта, обеспечивающий удобные поездки для пассажиров при минимизации операторских затрат \cite{mandl1980tndp,duranMicco2022tndpReview,sahin2023linePlanning,morozov2026oneRuleTndp}.

Эксперимент для арктической системы расселения основан на модели сети, заданной внешней средой (EFN), предназначенной для анализа сезонной устойчивости сервисной доступности в условиях климатически ограниченной транспортной связности \cite{mineEnvironmentFramedNetworks2026}. Результаты моделирования показывают, что системы расселения в рассматриваемых арктических регионах демонстрируют разные паттерны уязвимости. Нестабильность сети, вызванная сезонными изменениями температуры, приводит к изменениям отношений ``потребитель--поставщик''.

Рисунок \ref{fig:chapter4-tikhevich-ga-results} показывает, что за счет оптимизации времени доступности между кварталами требуемое число новых сервисов было уменьшено с 13 до 9. Таким образом, предложенные алгоритмом улучшения связей между кварталами следует интерпретировать как предложения по улучшению транспортной связности территории. Необходимо отметить, что лучшие решения, требующие минимального количества дополнительных сервисов включают в себя как только добавление маршрута, так и добавление маршрута с улучшением дорожной сети.

Эксперимент был проведен на 7 городах с разной плотностью застройки, населением и развитостью социальной инфраструктуры. Видно, что лишь в 1 из 7 случаев использование средне-взвешенной связности дало результат не хуже, чем использование сервис-ориентированной связности. Важно отметить, что результаты показывают внутреннюю структуру спроса, которая не всегда может быть одинаково распределена в пространстве.
